% ===========================================================
% Chapter 17: Open Problems and Emerging Directions
% Part IV: Frontiers
% ===========================================================

\chapter{Open Problems and Emerging Directions}
\label{ch:open}

\begin{quote}
\itshape
The best way to predict the future is to invent it.
\par\medskip
\upshape\raggedleft
---\,Alan Kay
\end{quote}

\bigskip

The preceding chapters have traced reinforcement learning from
multi-armed bandits (Chapter~\ref{ch:bandits}) through dynamic
programming (Chapter~\ref{ch:dp}), temporal-difference methods
(Chapter~\ref{ch:td}), deep value-function approximation
(Chapter~\ref{ch:value-approx}), policy gradients and actor-critic
algorithms (Chapters~\ref{ch:pg}--\ref{ch:actor-critic}), and onward
to the modern intersection of RL with large language models
(Chapters~\ref{ch:lm-post-training}--\ref{ch:agents}).  Along the
way we have seen RL solve games, align language models with human
preferences, enable complex reasoning, and coordinate multi-agent
systems.

Yet for every problem RL has solved, a dozen remain open.  Some are
long-standing\,---\,sample efficiency has been a concern since the
earliest tabular algorithms\,---\,while others have emerged only
recently, driven by the rapid adoption of RL as a post-training
method for foundation models.  This final chapter surveys the most
important open questions, groups them into coherent themes, and
attempts to sketch the road ahead.  It is necessarily more speculative
than the chapters that precede it: the goal is not to derive theorems
but to orient the reader toward the problems that will shape the next
decade of research.


% ===========================================================
\section{Sample Efficiency}
\label{sec:open-sample}
% ===========================================================

\index{sample efficiency}
Reinforcement learning is notoriously data-hungry.  Model-free
algorithms such as PPO (Chapter~\ref{ch:actor-critic})
and DQN (Chapter~\ref{ch:value-approx}) often require millions or
even billions of environment interactions to learn behaviours that a
human child acquires in minutes.  The gap between human and machine
sample efficiency is one of the oldest open problems in the field, and
it remains one of the most consequential.

\subsection{Model-Based RL and World Models}
\index{model-based reinforcement learning}
\index{world model}

The most natural route to sample efficiency is to \emph{learn a model
of the environment} and plan inside it.  If the agent can predict what
will happen before it acts, every real interaction becomes far more
informative.  Early model-based methods learned transition dynamics in
tabular or linear form (Chapter~\ref{ch:dp} discusses the planning
side); modern approaches learn \emph{world models}\,---\,neural
networks that predict future observations, rewards, and termination
signals.

Ha and Schmidhuber's ``World Models'' paper (2018) demonstrated that an
agent could learn a compressed latent representation of its
observations and then train a policy entirely inside the learned
model.  Hafner et al.\ extended this line of work with the Dreamer
family of algorithms (Dreamer, DreamerV2, DreamerV3), which learn a
recurrent state-space model and optimise a policy by
back-propagating through imagined trajectories in latent space.
DreamerV3 achieved human-level performance on dozens of continuous and
discrete control tasks\,---\,and even mined diamonds in Minecraft\,---\,using
roughly $10$--$100\times$ fewer environment steps than comparable
model-free methods.

\index{Dreamer}
\index{MuZero}
MuZero (Schrittwieser et al., 2020) took a different approach:
rather than predicting raw observations, it learned a latent dynamics
model whose hidden states were trained end-to-end to support planning
via Monte Carlo tree search.  MuZero matched the performance of
AlphaZero in Go, chess, and shogi while also mastering Atari games,
demonstrating that latent-space planning can rival or exceed
model-free methods even in visually complex domains.

Despite these successes, model-based RL faces persistent challenges.
Compounding model error\,---\,small inaccuracies in the learned
dynamics that accumulate over long planning horizons\,---\,can produce
policies that are optimal in the model but catastrophic in the real
environment.  Deciding \emph{when} to trust the model and when to
fall back on real data (the so-called Dyna trade-off; Sutton, 1991)
remains largely heuristic.

\subsection{Transfer and Few-Shot RL}
\index{transfer learning!in RL}
\index{few-shot RL}

A complementary route to sample efficiency is \emph{transfer}:
leveraging knowledge from previous tasks to accelerate learning on new
ones.  Meta-RL algorithms such as MAML (Finn et al., 2017) and
RL$^2$ (Duan et al., 2016) learn an initialisation or an inner-loop
learning algorithm that can adapt to a new task from a handful of
episodes.  More recently, foundation-model approaches treat RL
trajectories as sequences and use in-context learning\,---\,the same
mechanism that enables few-shot prompting in language models\,---\,to
adapt to new environments at inference time without gradient updates
(see Section~\ref{sec:open-offline} for the Decision Transformer and
related methods).

The vision of a single, general-purpose RL agent that can transfer
across diverse domains\,---\,an ``RL foundation model''\,---\,remains
largely unrealised.  Gato (Reed et al., 2022) showed that a single
transformer could play Atari, caption images, and control a robotic
arm, but its per-task performance lagged behind specialist agents.
Closing this gap is an active area of research.


% ===========================================================
\section{Reward Specification and Alignment}
\label{sec:open-alignment}
% ===========================================================

\index{reward hacking}
\index{Goodhart's law}
\index{alignment}
Every RL algorithm optimises a reward signal.  When that signal
faithfully captures what we want, optimisation is beneficial; when it
does not, the agent exploits the gap\,---\,a phenomenon known as
\emph{reward hacking}.  The problem is an instance of Goodhart's
law\index{Goodhart's law}: ``When a measure becomes a target, it
ceases to be a good measure.''

Chapter~\ref{ch:reward-models} introduced reward models trained from
human comparisons, and Chapter~\ref{ch:dpo} showed how DPO sidesteps
explicit reward modelling.  Both approaches mitigate reward hacking to
some extent, but neither eliminates it.  As models become more
capable, the stakes of misspecified rewards grow.

\subsection{Scalable Oversight}
\index{scalable oversight}

When the agent's capabilities exceed the evaluator's, human feedback
becomes unreliable.  A coding model may produce solutions that are
correct but inscrutable; a scientific-reasoning model may propose
experiments whose validity a non-expert cannot judge.  The
\emph{scalable oversight} problem asks: how can we provide reliable
training signal for systems whose outputs we cannot easily verify?

Several research programmes address this question.  \emph{Debate}
(Irving et al., 2018) pits two models against each other, with a
human judge deciding which argument is more convincing; the hope is
that the truth is easier to verify than to generate.
\emph{Recursive reward modelling} (Leike et al., 2018) uses an
AI assistant to help the human evaluate, creating a hierarchy of
oversight.  \emph{Constitutional AI} (Bai et al., 2022) replaces
some human feedback with a set of written principles that the model
critiques itself against, reducing\,---\,though not
eliminating\,---\,the burden on human evaluators.

\subsection{The Fundamental Challenge}

At its core, the alignment problem is one of \emph{value
specification}\index{value specification}: translating the rich,
context-dependent, and often contradictory landscape of human values
into a formal objective.  RL provides the optimisation machinery, but
it cannot tell us \emph{what} to optimise.  This is not merely a
technical problem; it is a philosophical one, touching on questions
that ethicists and political theorists have debated for centuries.

The practical implication for RL researchers is that reward engineering
will remain a first-class concern.  Advances in interpretability,
better methods for eliciting human preferences, and formal frameworks
for specifying constraints will all play a role.  No single technique
is likely to ``solve'' alignment; the challenge is systemic, and
progress will be incremental.


% ===========================================================
\section{Safety and Robustness}
\label{sec:open-safety}
% ===========================================================

\index{safe reinforcement learning}
\index{robustness}
An RL policy that performs well on average may fail catastrophically
in rare but important situations.  Safety and robustness are
therefore not luxuries\,---\,they are prerequisites for deploying RL
in the real world.

\subsection{Safe Exploration}
\index{safe exploration}
\index{constrained MDP}

In domains such as autonomous driving, healthcare, and industrial
control, the agent cannot be allowed to explore freely: certain states
(collisions, overdoses, equipment damage) must be avoided even during
training.  \emph{Constrained MDPs}\index{constrained MDP} formalise
this by augmenting the standard MDP with cost functions and
constraints of the form $\E_\pi[\sum_t c(s_t, a_t)] \leq d$.
Lagrangian methods and constrained policy optimisation (CPO;
Achiam et al., 2017) optimise the reward subject to these constraints,
but guaranteeing constraint satisfaction during learning\,---\,not
just at convergence\,---\,remains difficult.

\subsection{Distributional Robustness and Adversarial Attacks}
\index{distributional robustness}
\index{adversarial attack!on RL}

Policies trained in one environment may break when deployed in another
that differs even slightly.  \emph{Distributional robustness}
formulates the problem as a game between the agent and an adversary
that perturbs the transition dynamics or observations within a
bounded uncertainty set.  Robust MDPs (Iyengar, 2005; Nilim and
El~Ghaoui, 2005) provide a theoretical foundation, and recent work
has scaled these ideas to deep RL via adversarial training and
domain randomisation.

A related concern is \emph{adversarial attacks on RL policies}:
small, carefully crafted perturbations to the observation can cause a
well-trained policy to take catastrophically bad actions (Huang et
al., 2017).  Certified defences analogous to those developed for
image classifiers are still in their infancy for sequential
decision-making.

\subsection{Formal Verification}
\index{formal verification}

Can we \emph{prove} that a learned policy satisfies a safety property?
Formal verification techniques from software engineering\,---\,model
checking, abstract interpretation, satisfiability modulo theories
(SMT)\,---\,have been applied to small neural-network controllers, but
scaling them to the networks used in modern deep RL remains an open
problem.  Bridging the gap between the statistical guarantees of
learning theory and the deterministic guarantees of formal methods is
one of the most important\,---\,and most difficult\,---\,challenges in
the field.


% ===========================================================
\section{Offline and Batch RL}
\label{sec:open-offline}
% ===========================================================

\index{offline reinforcement learning}
\index{batch reinforcement learning}
In many practical settings\,---\,healthcare, recommendation systems,
dialogue\,---\,online interaction is expensive or risky.
\emph{Offline RL} (also called \emph{batch RL}) aims to learn a
good policy from a fixed dataset of previously collected transitions,
with no further environment interaction.  The idea is appealing:
organisations often possess large logs of past behaviour, and offline
RL promises to extract value from this data.

\subsection{The Distribution-Shift Problem}
\index{distribution shift!offline RL}

The central difficulty is \emph{distribution shift}.  The dataset was
collected by some behaviour policy $\mu$, but the learned policy
$\pi$ may visit states that $\mu$ rarely or never encountered.  In
those regions, the value estimates are unreliable, and naive off-policy
methods can diverge or produce delusionally optimistic policies.

\emph{Conservative Q-Learning} (CQL; Kumar et al., 2020) addresses
this by adding a regularisation term that penalises the Q-function on
out-of-distribution actions, effectively learning a lower bound on the
true value.  Implicit Q-Learning (IQL; Kostrikov et al., 2022)
avoids querying the policy altogether during training, using expectile
regression to estimate upper-envelope Q-values from the dataset.
Both methods have shown strong performance on standard benchmarks,
but their effectiveness depends heavily on the quality and coverage
of the dataset.

\subsection{Sequence Modelling and Decision Transformers}
\index{Decision Transformer}

An influential alternative to value-function-based offline RL is to
reframe the problem as \emph{sequence modelling}.  The Decision
Transformer (Chen et al., 2021) conditions a transformer on the
desired return, past states, and past actions, and autoregressively
predicts future actions.  At test time, the agent conditions on a
high return and generates the corresponding behaviour.  This approach
sidesteps Bellman backups entirely, leveraging the sequence-modelling
strengths of transformers.

Trajectory Transformer (Janner et al., 2021) extended the idea by
modelling the joint distribution over states, actions, and rewards,
enabling beam-search planning at inference time.  These methods blur
the line between RL and supervised learning, and their theoretical
properties\,---\,when they outperform, when they fail, and why\,---\,are
still being understood.

\subsection{Offline-to-Online}
\index{offline-to-online RL}

A promising hybrid is \emph{offline-to-online} RL: pre-train a policy
from a static dataset, then fine-tune with a small amount of online
interaction.  This mirrors the pre-train-then-fine-tune paradigm that
has been so successful for language models
(Chapter~\ref{ch:lm-post-training}).  Early results (Nair et al.,
2020; Lee et al., 2022) show that offline pre-training can
dramatically accelerate online learning, but preventing catastrophic
forgetting of offline knowledge and managing the distribution shift
during the transition remain open problems.


% ===========================================================
\section{Multimodal and Embodied RL}
\label{sec:open-embodied}
% ===========================================================

\index{embodied RL}
\index{sim-to-real transfer}
Reinforcement learning began with physical intuitions\,---\,a robot
navigating a maze, a thermostat regulating temperature\,---\,but for
the past decade the most visible successes have been in simulation:
games, language, and recommendation.  Returning RL to the physical
world is one of the field's great unfinished projects.

\subsection{Sim-to-Real Transfer}

Training in simulation is cheap and safe; deploying in the real world
is neither.  The \emph{sim-to-real gap}\index{sim-to-real gap}
arises because simulators inevitably simplify physics, sensor noise,
and visual appearance.  Domain randomisation (Tobin et al., 2017)
mitigates the gap by training the policy across a distribution of
simulated environments, so that the real world appears as just another
sample.  System identification, adaptive control, and learned residual
models offer complementary approaches, but reliably closing the
sim-to-real gap for contact-rich manipulation remains an active
challenge.

\subsection{Vision-Language-Action Models}
\index{vision-language-action model}

A new class of \emph{vision-language-action} (VLA) models aims to
unify perception, language understanding, and motor control in a
single architecture.  RT-2 (Brohan et al., 2023) fine-tuned a
vision-language model to output robotic actions as text tokens,
demonstrating that semantic knowledge from web-scale pre-training
could transfer to manipulation tasks.  Octo (Ghosh et al., 2024) and
$\pi_0$ (Black et al., 2024) pushed this further with
transformer architectures that consume images and language instructions
and emit continuous joint commands.

These models are still far from the dexterity and adaptability of
human manipulation.  Key bottlenecks include the scarcity of
high-quality robotic demonstration data, the difficulty of specifying
rewards for contact-rich tasks, and the computational cost of
real-time inference on resource-constrained hardware.

\subsection{Continuous Control at Scale}

Beyond manipulation, RL is being applied to locomotion (legged
robots trained via PPO in simulation and deployed on hardware),
autonomous racing (Gran Turismo), drone navigation, and whole-body
humanoid control.  The common thread is that these tasks involve
high-dimensional continuous action spaces, complex contact dynamics,
and stringent real-time requirements.  Scaling RL to these settings
pushes both algorithmic and engineering frontiers.


% ===========================================================
\section{RL and Foundation Models}
\label{sec:open-foundation}
% ===========================================================

\index{foundation model!RL for}
Perhaps the most consequential development in recent RL is its
emergence as the \emph{universal post-training method} for foundation
models.  Chapters~\ref{ch:lm-post-training}--\ref{ch:reasoning}
documented this for text; the same pattern is now repeating across
modalities.

\subsection{RL for Language and Reasoning}

RLHF (Chapter~\ref{ch:dpo}) and GRPO (Chapter~\ref{ch:grpo}) have
become standard tools for aligning and improving large language models.
As discussed in Chapter~\ref{ch:reasoning}, RL with verifiable rewards
has enabled dramatic improvements in mathematical reasoning and code
generation.  The frontier is moving toward longer-horizon reasoning
tasks\,---\,multi-step scientific problem-solving, software
engineering, theorem proving\,---\,where the reward is sparse and
delayed, and the action space is the space of all possible text
continuations.

\subsection{RL for Visual Generation}

The same RLHF loop that improved language models has been applied to
image generation.  DPOK (Fan et al., 2023) and DDPO (Black et al.,
2024) used policy-gradient methods to fine-tune diffusion models
against human-preference reward models, improving aesthetic quality and
prompt adherence.  For video generation, RL-based fine-tuning has been
used to improve temporal coherence and visual fidelity, though the
computational cost of generating and scoring video trajectories remains
a bottleneck.

\subsection{RL for Scientific Discovery}
\index{scientific discovery!RL for}

RL has contributed to notable scientific breakthroughs.  AlphaFold's
(Jumper et al., 2021) training pipeline incorporated RL-like
structural refinement.  GNoME (Merchant et al., 2023) used RL-guided
search to discover millions of new stable crystal structures.
FunSearch (Romera-Paredes et al., 2024) combined large language models
with evolutionary search to discover new mathematical constructions,
outperforming human mathematicians on the cap-set problem.

The common pattern is RL as the \emph{outer loop}: a foundation model
proposes candidates (protein structures, material compositions,
mathematical conjectures), and an RL-like objective scores and selects
among them.  Extending this pattern to chemistry, drug discovery,
climate modelling, and engineering design is an active and promising
frontier.


% ===========================================================
\section{Theoretical Frontiers}
\label{sec:open-theory}
% ===========================================================

\index{RL theory}
The theory of RL is rich in the tabular and linear-function-approximation
settings (Chapters~\ref{ch:mdp}--\ref{ch:td}), but the methods that
work best in practice\,---\,deep networks, transformers, high-dimensional
policy spaces\,---\,often outrun available theory.

\subsection{Regret and Sample Complexity for Deep RL}
\index{regret bounds!deep RL}
\index{PAC-MDP}

Regret bounds for tabular MDPs are tight: algorithms such as UCBVI
achieve $\tilde{O}(\sqrt{H^3 S A T})$ regret, where $H$ is the
horizon, $S$ the number of states, $A$ the number of actions, and $T$
the number of episodes.  For linear MDPs, matching upper and lower
bounds are also known.  But for the nonlinear function approximation
used in deep RL\,---\,overparameterised neural networks trained with
stochastic gradient descent\,---\,no comparable guarantees exist.

\emph{Neural tangent kernel} (NTK) analyses provide some insight by
showing that sufficiently wide networks behave approximately linearly,
but this regime does not capture the representation learning that
makes deep RL effective.  Developing a theory that accounts for
feature learning, non-convex optimisation, and the
exploration--exploitation trade-off simultaneously is one of the
hardest open problems in machine learning theory.

\subsection{Convergence of Nonlinear Function Approximation}
\index{deadly triad}

The \emph{deadly triad}\,---\,the combination of bootstrapping,
off-policy learning, and function approximation\,---\,was identified
by Sutton and Barto as a source of divergence in approximate dynamic
programming (Chapter~\ref{ch:value-approx}).  While algorithms such
as DQN work well in practice, proving their convergence under
realistic assumptions (finite-width networks, non-i.i.d.\ data,
experience replay) remains elusive.  Recent progress includes
convergence results for overparameterised networks in the NTK regime
and two-timescale analyses of actor-critic methods, but a unified
theory is still missing.

\subsection{Statistical Complexity of RLHF}
\index{RLHF!statistical complexity}

The explosion of RLHF in language-model training
(Chapter~\ref{ch:dpo}) has spurred new theoretical questions.  How
many human comparisons are needed to learn a reward model that
faithfully represents human preferences?  Under what conditions does
the DPO objective recover the same policy as the RLHF pipeline?  How
does the quality of the preference data affect the statistical
efficiency of alignment?  Early results (Zhu et al., 2024; Xu et al.,
2024) establish sample-complexity bounds under the
Bradley--Terry\index{Bradley--Terry model} preference model, but the
theory does not yet account for the distributional and strategic
complexities of real human feedback.

\subsection{Why Does RL Work for LLMs?}

Perhaps the deepest theoretical puzzle is why RL fine-tuning works as
well as it does for large language models.  The action space is
enormous (the full vocabulary at each token position), the reward is
sparse and delayed (often a single scalar for an entire response), and
the policy is a billion-parameter network.  Classical RL theory
predicts that such problems should be intractable, yet PPO and GRPO
produce reliable improvements in practice.  Understanding
this\,---\,likely through a combination of the structure of natural
language, the implicit exploration provided by token sampling, and the
warm start from pre-training\,---\,is a key open question.


% ===========================================================
\section{The Road Ahead}
\label{sec:open-road}
% ===========================================================

Let us close by stepping back from individual problems and considering
the broader trajectory of the field.

\paragraph{RL as universal fine-tuning.}
The most striking trend of the past several years is the
generalisation of RL from a niche method for game-playing and robotics
into a \emph{universal fine-tuning method} for any model that generates
structured output.  Language models, image generators, video
synthesizers, protein-structure predictors, and robotic controllers are
all being improved by RL-based post-training.  The underlying
principle is simple: if you can score the output, you can use RL to
make the model produce higher-scoring outputs.  This principle, first
formalised in the policy-gradient theorem (Chapter~\ref{ch:pg}), is
proving to be one of the most powerful ideas in modern AI.

\paragraph{Convergence of classical theory and modern practice.}
For decades, RL theory and RL practice developed in largely separate
communities.  The theorists studied tabular MDPs and linear function
approximation; the practitioners trained deep networks at scale.
This gap is beginning to close.  Offline RL theory informs practical
algorithm design; PAC-Bayes bounds provide non-vacuous generalisation
guarantees for learned policies; and the study of RLHF is bringing
learning-theoretic tools to bear on one of the most consequential
applications of RL.  A unified theory that connects the elegant
results of the tabular setting to the empirical successes of deep RL
would be a landmark achievement.

\paragraph{Agents and autonomy.}
Chapter~\ref{ch:reasoning} and Chapter~\ref{ch:agents} described
the emergence of RL-trained agents that can reason, use tools, and
coordinate with other agents.  This trend is accelerating.  As
language models become more capable and as RL fine-tuning enables
longer-horizon planning, we may see agents that autonomously conduct
scientific research, write and debug complex software, manage
infrastructure, and negotiate on behalf of their users.  The
technical challenges are immense\,---\,credit assignment over
thousands of steps, safe exploration in open-ended environments,
robust multi-agent coordination\,---\,but so are the potential
benefits.

\paragraph{The importance of safety.}
With greater autonomy comes greater risk.  The safety and alignment
challenges discussed in Sections~\ref{sec:open-alignment}
and~\ref{sec:open-safety} will only become more pressing as agents
become more capable.  The field needs not only better algorithms but
also better institutions: evaluation standards, red-teaming
practices, regulatory frameworks, and a culture of responsible
development.  RL researchers have a special responsibility here, because
the optimisation pressure that RL exerts is among the strongest in
all of machine learning.

\paragraph{What the next decade might bring.}
Predicting the future of a rapidly moving field is a fool's errand,
but a few trends seem likely to persist.  First, the scale of RL
training will continue to grow, driven by the economics of
post-training foundation models.  Second, the boundary between RL,
supervised learning, and self-supervised learning will continue to
blur, as methods like DPO (Chapter~\ref{ch:dpo}) and Decision
Transformer (Section~\ref{sec:open-offline}) demonstrate.  Third, RL
will be applied to an ever-wider range of domains, from materials
science to climate policy to education.  Fourth, the safety and
alignment challenges will force the community to develop new tools for
evaluation, oversight, and governance.

The story of reinforcement learning is far from over.  If anything, it
is just reaching its most interesting chapter.  The mathematical
foundations laid in Parts~I and~II of this book, the algorithmic
toolkit developed in Part~III, and the frontier applications surveyed
in Part~IV all provide the vocabulary needed to contribute.  The open
problems described in this chapter are not obstacles to be lamented;
they are opportunities to be seized.  We hope this book has equipped
you to seize them.
